\scanpage{79}
pt de ramification d'indice $p$ un diviseur de degré $2$ stable par $G_0$, encore absurde. Donc on est dans l'un des trois cas, $(2,3,3),(2,3,4),(2,3,5)$.

II) \deleted{\unsure{Considérons} \unsure{Au} \unsure{voisin}}\inserted{Nous allons} montrer que ces trois cas \inserted{de ss-groupes finis de $\operatorname{GP}(1,\mathbb C)$} correspondent à des classes de conjugaison bien déterminées de ss-groupes de $\operatorname{GP}(1,k)$ : groupes tétraédral, octahédral, icosaédral.

Pour ceci, considérons les couples $(\nu_0,\nu_1)$ correspondants (cas I, II, III) \deleted{\unsure{en|---}} choisissant dans le cas III $\alpha=\zeta+\zeta^{-1}$, où $\zeta=\exp\frac{2i\pi}{5}$, donc $\alpha=2\cos\frac{2\pi}{5}$. Les \deleted{\unsure{vérifions que les}}\inserted{\unsure{supposons connus|supposés connus}} les polyèdres réguliers \unsure{combinatoires|correspondants} \unsure{s'identifiant|s'identifient} aux tétraèdre, octaèdre, icosaèdre -- on constate aisément, sur les formules explicites qu'on trouve des réalisations géom \inserted{strictement} fidèles de ceux-ci. Leurs groupes \inserted{$G_0$} sont respectivement

$\unsure{\mathfrak S_4|\mathscr S_4},\ \mathfrak S_4\times\mathbb Z/2\mathbb Z,\ \mathfrak A_5\times\mathbb Z/2\mathbb Z$, et \deleted{\unsure{l'}}\inserted{les} images de celui-ci dans $\operatorname{GP}(1,k)$ \deleted{est} sont respectivement $\simeq\unsure{\mathfrak A_4|\mathfrak S_4},\mathfrak S_4,\mathfrak A_5$

Appelons ceux-ci les groupes pythagoriciens standart dans $\operatorname{GP}(1,k)$ (définis mod conjugaison). Ceci a établi \inserted{en passant} l'existence des cas I, II, III, et par dualité II') III') -- \deleted{l'unicité}

Il reste à prouver que le seul cas de pol. régulier fini qui n'est pas couvert par les précédents est le cas IV -- et d'en prouver l'existence de IV. Or \inserted{pour \unsure{celle-ci|celui-ci},} la construction du pentagramme associé à un icosaèdre est bien claire. -- cf plus bas la forme que prend cette construction dans notre contexte actuel

\deleted{III) Revenons au cas de $G_0\subset\operatorname{GP}$ \unsure{, engendré par|, engendré pour} $\tau_0,\tau_1,\tau_2$, fini.}
% archive sheet number: 78
